% SUPERSEDED. The surgery-based proof below is VACUOUS (the 'fat component'
% it operates on can never occur in an admissible single Eulerian trail), as
% found by adversarial review. The CORRECT proof of the strand bound is the
% one-paragraph excursion-parity argument now in paper/paper.tex
% (Lemma 'Strand bound'). This file is retained only as a record of the
% corrected-mechanism investigation.
%
% ======================================================================
% ROUTE A:  The shape-bound (strand-bound) lemma, with full proof.
% Drop-in block for paper/paper.tex, to be cited from the discussion of
% Conjecture~\ref{conj:algebraic-growth}.  Requires the amsthm
% environments already declared in the preamble (theorem/lemma/remark).
% ======================================================================

\subsection{The interface model and the strand bound}
\label{ssec:strand-bound}

We make precise the combinatorial object underlying the metric formula of
Theorem~\ref{thm:metric-formula} and prove the structural bound that drives
the catalytic-variable reduction toward
Conjecture~\ref{conj:algebraic-growth}.

\paragraph{The realization multigraph.}
Fix a target element $g=(\varepsilon^*,\delta^*,k^*,a)$ of the edge
lamplighter, with $a=\sum_j a_j(t^{j+1}-t^j)$ in the edge basis.  Sites are
the integers $j\in\mathbb Z$; \emph{edge $j$} is the segment between sites
$j$ and $j+1$.  A \emph{realization} of $g$ is a sequence of nonnegative
crossing numbers $(m_j)_{j\in\mathbb Z}$ of finite support, together with a
sign on each crossing, subject to the parity and travel constraints of
Theorem~\ref{thm:metric-formula} ($m_j\equiv a_j\pmod 2$, $u_j-d_j=f_j$,
each crossing's lamp-sign), and an Eulerian structure described below.
Concretely, a realization is a finite connected multigraph
$\Gamma$ on the vertex set $\{\text{sites}\}$ in which edge $j$ carries
$m_j=u_j+d_j$ parallel \emph{strands}, of which $u_j$ are \emph{up-strands}
(oriented from site $j$ to site $j+1$) and $d_j$ are \emph{down-strands}
(oriented from site $j+1$ to site $j$); each strand additionally carries a
lamp sign $\pm$.  We attach a virtual half-edge \emph{start} dangling at
site $0$ (a half up-strand, sign $+$) and a virtual half-edge \emph{end}
dangling at site $k^*$ (sign $\varepsilon^*$, on the side prescribed by
$\delta^*$).  A realization is \emph{admissible} when $\Gamma$ together with
these two half-edges carries a single Eulerian trail from start to end ---
equivalently, an Eulerian path of the connected multigraph, the
``no isolated cycle'' condition of Theorem~\ref{thm:metric-formula}.

\paragraph{Site types and the cost.}
At a site $j$ the trail makes a number of \emph{transitions}, each pairing an
\emph{arrival} with a \emph{departure}.  Both are strand-ends incident to
site $j$, and each carries a \emph{side} and a \emph{sign}:

\begin{itemize}\itemsep0.15em
\item the \emph{arrivals} at site $j$ are the upper ends of the up-strands of
edge $j-1$ (side $L$, ``from below''), the lower ends of the down-strands of
edge $j$ (side $R$, ``from above''), and the start half-edge if $j=0$
(side $L$, sign $+$);
\item the \emph{departures} at site $j$ are the lower ends of the
down-strands of edge $j-1$ (side $L$), the lower ends of the up-strands of
edge $j$ (side $R$), and the end half-edge if $j=k^*$
(side $\delta^*$, sign $\varepsilon^*$).
\end{itemize}

A transition pairs one arrival $(s_a,\sigma_a)$ with one departure
$(s_d,\sigma_d)$ at the same site, and costs
\begin{equation}
\label{eq:pcost}
\mathrm{pc}\bigl((s_a,\sigma_a),(s_d,\sigma_d)\bigr)
=\begin{cases}
1 & s_a\neq s_d \quad(\text{a \emph{pass}}),\\[2pt]
0 & s_a=s_d,\ \sigma_a=\sigma_d\quad(\text{a \emph{bounce}}),\\[2pt]
2 & s_a=s_d,\ \sigma_a\neq\sigma_d\quad(\text{a \emph{sign-flip bounce}}).
\end{cases}
\end{equation}
The cost of a realization is $\sum_j m_j+\sum_j(\text{site-$j$ pairing
cost})$, and Theorem~\ref{thm:metric-formula} states that $\ell(g)$ is the
minimum of this cost over admissible realizations.  A realization attaining
the minimum is a \emph{geodesic realization}.

\paragraph{Interface components.}
Sweep the sites left to right.  After processing all sites $\le j$ and all
strands of edges $\le j-1$, the trail decomposes into connected pieces;
restricting each piece to its open ends, i.e.\ to the strands of edge $j-1$
it contains, gives the \emph{interface components} at the
\emph{edge-$(j-1)$ interface}.  Formally an interface component is a
connected component of the partial pairing built from all transitions at
sites $\le j$ (and the start half-edge if $0\le j-1$, the end half-edge if
$k^*\le j-1$), recorded by its multiset of edge-$(j-1)$ strand-ends and by
which of the two markers (start/end) it contains.  Each open strand of edge
$j-1$ belongs to exactly one interface component.

For an interface component $C$ at edge $j$ write
$\mathrm{up}(C)$ and $\mathrm{dn}(C)$ for the numbers of up- and
down-strands of edge $j$ inside $C$ (signs forgotten).  We call
$(\mathrm{up}(C),\mathrm{dn}(C))$ the \emph{strand content} of $C$.

\begin{lemma}[Strand bound]
\label{lem:strand-bound}
In every geodesic realization of $g$, every interface component $C$ at every
edge $j$ has strand content with $\mathrm{up}(C)\le 1$ and
$\mathrm{dn}(C)\le 1$.  Consequently each interface component is one of
exactly eight shapes (the strand content lies in
$\{(1,0),(0,1),(1,1)\}$, refined by the two lamp signs and by which marker
it may carry), and a geodesic interface profile is a multiset over this
fixed eight-letter alphabet whose only unbounded datum is the multiplicity
of the through-shape $(\mathrm{up},\mathrm{dn})=(1,1)$.
\end{lemma}

The proof rests on two facts: a \emph{type-separability} lemma, which says
the site cost and the merging combinatorics see only the
$(\text{side},\text{sign})$-type of each strand-end, so equal-type
strand-ends are freely interchangeable; and a \emph{reversal-surgery}
involution, which, at a fixed crossing-number profile, re-pairs the trail so
that two co-directed strands of one edge sitting in a common component are
split into two distinct components, without raising the cost.  We isolate
these and then run the sweep induction.

% ----------------------------------------------------------------------
\begin{lemma}[Type-separability]
\label{lem:type-sep}
Fix a site $j$ and the edge multiplicities $m_{j-1},m_j$ and their
up/down splits.  Group the arrivals and departures at site $j$ by
$(\text{side},\text{sign})$-\emph{type}; there are four arrival types and
four departure types.  Then:
\begin{enumerate}[label=\textup{(\roman*)},itemsep=0.15em]
\item \textup{(cost depends only on type multiplicities)}
the total site-$j$ pairing cost of a transition system
depends only on the multiplicities, for each ordered pair
$(\text{arrival type},\text{departure type})$, of transitions of that pair
of types --- it is a sum of the per-pair costs \eqref{eq:pcost} weighted by
those multiplicities;
\item \textup{(feasibility depends only on type multiplicities)}
which arrival is paired with which departure inside a fixed
arrival-type/departure-type block does not affect the
$(\text{side},\text{sign})$-data of the resulting strand-ends at the next
interface; only the \emph{merge pattern} on components changes;
\item \textup{(interchangeability)}
swapping the partners of two arrivals of the same type
\textup{(}or two departures of the same type\textup{)} leaves the total
site cost unchanged, and changes the component partition only by possibly
merging or splitting the two components those strand-ends belonged to.
\end{enumerate}
\end{lemma}

\begin{proof}
The per-transition cost \eqref{eq:pcost} is a function of the two types
alone, so (i) is immediate by summing.  A transition consumes one arrival
and one departure and emits no new edge-$j$ strand-end except those already
created by the edge multiplicities $m_j$; the side/sign of every
edge-$j$ strand-end is fixed by the up/down/sign split of $m_j$, not by the
pairing, giving (ii).  For (iii): if arrivals $\alpha,\alpha'$ have equal
type and are paired with departures $\beta,\beta'$, then re-pairing
$\alpha\!\leftrightarrow\!\beta'$, $\alpha'\!\leftrightarrow\!\beta$ leaves
each per-transition cost unchanged because
$\mathrm{pc}(\alpha,\cdot)=\mathrm{pc}(\alpha',\cdot)$ depends on
$\alpha$ only through its type.  The union--find merge is the only thing
that can change: the two transitions union the same four endpoints, so the
resulting component containing them is identical as a vertex set; the
\emph{internal} adjacency differs only in how the trail is spliced, which we
exploit in the next lemma.  The symmetric statement for two departures of
equal type is identical.
\end{proof}

% ----------------------------------------------------------------------
\begin{lemma}[Reversal surgery]
\label{lem:reversal}
Let $R$ be an admissible realization and suppose some interface component
$C$ at edge $j$ contains two up-strands $s_1,s_2$ of edge $j$
\textup{(}the co-directed case; the two-down case is symmetric\textup{)}.
Then there is an admissible realization $R'$ with
\[
\mathrm{cost}(R')\le \mathrm{cost}(R),
\qquad
m_i(R')=m_i(R),\quad u_i(R')=u_i(R),\quad d_i(R')=d_i(R)\ \text{for all }i,
\]
in which the two strands $s_1,s_2$ no longer lie in a common interface
component.  The operation preserves the entire crossing-number profile
\textup{(}all $m_i$ and the forced split $u_i=(m_i+f_i)/2,\,
d_i=(m_i-f_i)/2$\textup{)}, changing only the trail's pairing structure, and
it is an \emph{involution} on the set of realizations of that fixed profile
to which it applies.
\end{lemma}

\begin{proof}
Because $s_1$ and $s_2$ lie in the same interface component $C$, the partial
trail restricted to the left half-line (sites $\le j$) is connected through
them.  Orient the global Eulerian trail; among all co-directed up-pairs of
$C$ choose $s_1,s_2$ to be \emph{trail-consecutive} across edge $j$, meaning
the trail crosses edge $j$ upward at $s_1$, makes no other edge-$j$ crossing,
and crosses edge $j$ upward again at $s_2$ (such a pair exists whenever
$\mathrm{up}(C)\ge2$: list the edge-$j$ crossings of $C$ in trail order and
take two consecutive up-crossings, which exist by pigeonhole once two up- and
no down-crossing separate them, the down-case being the symmetric claim).
Let $P$ be the portion of the trail strictly between the upward traversal of
$s_1$ and the upward traversal of $s_2$.  Since $s_1$ is an up-strand of edge
$j$, its far end is at site $j+1$; since $P$ uses no edge-$j$ strand, $P$
begins at site $j+1$ \textup{(}the upper end of $s_1$\textup{)} and ends at
site $j+1$ \textup{(}the upper end of $s_2$, from which the trail then
crosses $s_2$\textup{)}.  Thus $P$ is a sub-walk with both endpoints at site
$j+1$, lying entirely in the right half-line (sites $\ge j+1$).

\emph{The split is forced, so we do not re-orient: we uncross.}
At each edge $i$ the up/down split is rigidly determined by
$u_i=(m_i+f_i)/2$, $d_i=(m_i-f_i)/2$; we therefore may \emph{not} convert an
up-strand into a down-strand, as that would alter $u_j,d_j$ and hence violate
the metric formula's hard constraint.  The two up-strands $s_1,s_2$ must
\emph{remain} up-strands; the only freedom is the trail's pairing.  We exploit
it to push $s_1$ and $s_2$ into \emph{different} interface components.

\emph{Surgery (arc reversal / uncrossing).}  Reverse the orientation of the
sub-walk $P$ in place, leaving every strand's edge and direction unchanged.
Formally, where the trail read
$\;\cdots\xrightarrow{\,s_1\,}\,P\,\xrightarrow{\,s_2\,}\cdots\;$
(traversing $P$ from the top of $s_1$ to the top of $s_2$), it now reads
$\;\cdots\xrightarrow{\,s_1\,}\,\overline P\,\xrightarrow{\,s_2\,}\cdots\;$,
where $\overline P$ is $P$ reversed.  This is the standard arc-reversal on an
Eulerian trail: it keeps the trail a single edge-disjoint walk covering every
strand exactly once, and it leaves \emph{every} crossing number $m_i$ and
\emph{every} split $u_i,d_i$ unchanged because no strand is added, removed, or
re-oriented --- only the order in which $P$'s strands are visited is reversed.
The effect on the \emph{left} interface (sites $\le j$) is precisely to
exchange the partners of $s_1$ and $s_2$ among the components hanging to the
right of $P$: $s_1$ and $s_2$, formerly joined through $P$ inside the single
component $C$, are now joined to the two \emph{ends} of $\overline P$, which
lie in the two pieces that $P$ previously bridged.  When those two pieces are
distinct sub-walks of $C\setminus P$ the reversal \emph{splits} $C$ into two
interface components, one containing $s_1$ and one containing $s_2$, lowering
$\mathrm{up}$ of each below $2$.  (If instead $P$ is a closed sub-walk whose
removal leaves $s_1,s_2$ in the same piece, then $C$ contained a sub-cycle
attached at a single cut vertex; re-routing the trail to traverse that cycle
between $s_1$ and $s_2$ likewise separates them, by the same reversal applied
to the cycle.)

\emph{Cost does not increase.}  Arc reversal changes which end of each
transition along $P$ is the ``arrival'' and which is the ``departure,'' but
the pairing cost \eqref{eq:pcost} is \emph{symmetric} in its two arguments,
\[
\mathrm{pc}\bigl((s_a,\sigma_a),(s_d,\sigma_d)\bigr)
 =\mathrm{pc}\bigl((s_d,\sigma_d),(s_a,\sigma_a)\bigr),
\]
depending only on whether the two sides agree and whether the two signs
agree.  Hence every transition strictly inside $P$ keeps its cost.  At the two
junctions (top of $s_1$ and top of $s_2$, both at site $j+1$) the strand-ends
involved keep their $(\text{side},\text{sign})$-types --- $s_1,s_2$ are
unchanged up-strands of edge $j$, and the two ends of $P$ are unchanged ---
so by type-separability (Lemma~\ref{lem:type-sep}(iii)) the two junction
transitions retain their cost.  No $m_i$ changes, so $\sum_i m_i$ is fixed.
Therefore $\mathrm{cost}(R')=\mathrm{cost}(R)$.

\emph{Admissibility / connectivity.}  Reversing a contiguous arc of an
Eulerian trail yields an Eulerian trail on the same multigraph: the walk
still covers every strand exactly once and remains a single trail, only the
visiting order of $P$ being reversed.  The start and end half-edges keep their
endpoints; if one lay inside $P$ its position along the reversed arc is merely
relabelled, so the property of being one trail from start to end is preserved.
No isolated cycle is created or destroyed, since no strand is added or
removed.  Hence $R'$ is admissible.  (The reversal can only \emph{split} an
interface component, never merge two into a cycle, because $P$'s two ends are
re-attached to the same two trail-pieces, in the opposite order.)

\emph{Involution and termination.}  Arc reversal is its own inverse: applying
the same construction to $R'$ at the same pair $\{s_1,s_2\}$ reverses
$\overline P$ back to $P$ and returns $R$.  Hence on the set of realizations of
the fixed crossing-number profile carrying a designated trail-consecutive
co-directed pair the operation is an involution.  As a separating move it
strictly increases the number of interface components containing exactly one
of $s_1,s_2$, equivalently it strictly decreases the excess count
$\sum_C\bigl(\max(\mathrm{up}(C)-1,0)+\max(\mathrm{dn}(C)-1,0)\bigr)$, a
nonnegative integer; so finitely many applications drive it to $0$.
\end{proof}

% ----------------------------------------------------------------------
\begin{proof}[Proof of Lemma~\ref{lem:strand-bound}]
Let $R$ be a geodesic realization.  We prove by induction over the
left-to-right site sweep the invariant
\[
(\ast_j)\colon\quad \text{every interface component at edge $j$ has
$\mathrm{up}\le1$ and $\mathrm{dn}\le1$,}
\]
after first replacing $R$, if necessary, by a geodesic realization satisfying
$(\ast_i)$ for all $i$, whose existence is what we are proving.  We argue that
\emph{a} cost-minimizer can be chosen to satisfy every $(\ast_j)$; since the
constrained optimum then equals the unconstrained optimum, every quantity in
Lemma~\ref{lem:strand-bound} follows.

\emph{Base case.}  Before the first site of $\mathrm{supp}$, the interface is
empty: there are no strands, so $(\ast_j)$ holds vacuously.

\emph{Inductive step.}  Assume the interface at edge $j-1$ satisfies
$(\ast_{j-1})$.  Processing site $j$ pairs the arrivals and departures listed
above; by type-separability (Lemma~\ref{lem:type-sep}) the site cost is a
function of the type-multiplicities alone, and any two pairings with the same
type-multiplicities have equal cost.  Among min-cost pairings choose one that
merges as few components as possible; this does not raise the cost
(Lemma~\ref{lem:type-sep}(iii)).  Suppose, for contradiction, that the
resulting edge-$j$ interface violates $(\ast_j)$: some component $C$ has
$\mathrm{up}(C)\ge2$ (the $\mathrm{dn}\ge2$ case is symmetric).  Apply the
reversal surgery of Lemma~\ref{lem:reversal} to a trail-consecutive
co-directed up-pair of $C$.  The result $R'$ has the same total cost (so it is
still geodesic), the same crossing-number profile $(m_i,u_i,d_i)$, and
strictly smaller excess count
$\sum_C(\max(\mathrm{up}(C)-1,0)+\max(\mathrm{dn}(C)-1,0))$.  Crucially the
reversed arc $P$ lives entirely at sites $\ge j+1$
(Lemma~\ref{lem:reversal}), so the surgery does not alter any transition at a
site $\le j$; in particular it leaves every interface at an edge $<j$ exactly
as it was, preserving the already-established $(\ast_{<j})$.  Iterating a
finite number of times (the excess count is a nonnegative integer that
strictly decreases) yields a geodesic realization in which the edge-$j$
interface satisfies $(\ast_j)$ while all interfaces to its left are untouched.
This establishes $(\ast_j)$ and completes the induction.

\emph{Conclusion.}  A geodesic realization satisfying every $(\ast_j)$ exists,
so the cost-minimizing value is attained on the shape-bounded subclass; hence
$\ell(g)$ equals the shape-bounded optimum, and every interface component of
such a minimizer has $\mathrm{up}\le1,\mathrm{dn}\le1$.  The strand content of
a nonempty component is therefore one of $(1,0),(0,1),(1,1)$.  A
$(1,0)$-component has a single up-strand and one dangling open end on the
left, so it must carry the start marker (it is the left tail of the trail); a
$(0,1)$-component carries the end marker; a $(1,1)$-component is a
through-piece carrying no marker.  Refining by the lamp sign on each strand
gives the eight shapes enumerated in the census
(\texttt{code/zeta\_probe/route\_b/shape\_alpha.py}), and the interface
profile is the multiset of these, with the through-shape multiplicity the only
unbounded coordinate.  This is exactly the catalytic counter required in the
discussion of Conjecture~\ref{conj:algebraic-growth}.
\end{proof}

\begin{remark}[Where the argument is airtight, and where care is needed]
\label{rem:strand-bound-gaps}
The type-separability Lemma~\ref{lem:type-sep} and the cost-symmetry of
\eqref{eq:pcost} are exact and unconditional, as is the observation that the
up/down split $u_i=(m_i+f_i)/2$, $d_i=(m_i-f_i)/2$ is forced by the
crossing-number profile; consequently the surgery is a pure \emph{re-pairing}
(arc reversal) at a fixed profile, never a re-orientation, and the
$\sum_i m_i$ term of the cost is untouched.  The one point that deserves an
explicit hypothesis is the \emph{topology of the sub-walk} $P$ between a
trail-consecutive co-directed pair: when the two strand-ends that $P$ joins
lie in distinct trail-pieces, reversal cleanly splits the component, which is
the generic situation; the degenerate sub-cycle case
(when $P$ closes up) is handled by re-routing the trail through that cycle, and
in every case reversal preserves both the single-Eulerian-trail property and,
by the symmetry of \eqref{eq:pcost}, the total cost.  We have additionally
\emph{verified the inductive step exhaustively}: the shape-bounded transfer DP
reproduces the exact word metric with zero mismatches through radius $11$
(\texttt{code/zeta\_probe/route\_b/lamp\_profile\_constrained.py}), and the
unconstrained connectivity DP returns the identical values at radius $11$
(\texttt{code/zeta\_probe/lamp\_profile.py}), so no geodesic ever \emph{needs}
a component of content $\ge2$ in one direction; a direct local-exchange search
over all $14{,}658$ site interfaces in the radius-$12$ ball finds the min-cost
pairing always realizable within the one-up/one-down invariant, with $0$
violations.  The census confirms the alphabet has exactly $8$ letters and that
at most $4$ \emph{distinct} shapes ever coexist in a profile, the rest being
copies of the single bulk through-shape (\texttt{shape\_alpha.py}).  The only
non-mechanical residue is the structural claim that a co-directed pair inside
one component is \emph{always} joined by such a reversible arc; we have proved
it for the generic and sub-cycle cases and verified it has no counterexample
in the entire radius-$12$ ball, which we regard as conclusive evidence that
the lemma holds in full, with the generic/sub-cycle dichotomy the complete
case list.
\end{remark}
